\section{Crowd Counting}
% Crowd counting has emerged as one of the most significant and challenging tasks in computer vision, with applications spanning from public safety and security surveillance to urban planning and event management. Estimation of crowd count is becoming crucial nowadays, as it can help in security surveillance, crowd monitoring, and management for different events \cite{SALEH2015103}. The ability to accurately estimate crowd density and population size from visual data has become increasingly important in our interconnected world, where understanding human crowd dynamics is essential for maintaining public safety and optimizing resource allocation.\\
% The fundamental challenge in crowd counting lies in accurately estimating the number of individuals present in a single image or video frame, particularly in densely populated scenarios where traditional detection-based methods fail due to severe occlusions and overlapping figures. The task of crowd counting and density map estimation is riddled with many challenges such as occlusions, non-uniform density, intra-scene and inter-scene variations in scale and perspective \cite{SALEH2015103}. These challenges are further compounded by occlusion, scale variation, perspective distortion, rotation, illumination variation, and weather changes \cite{gao2025survey}, making crowd counting one of the most demanding tasks in computer vision.\\
% Traditional approaches to crowd counting have evolved through three main paradigms: detection-based methods, regression-based methods, and density estimation approaches. Early detection-based methods attempted to identify and count individual persons but proved inadequate for dense crowds due to severe occlusions. Regression-based methods sought to establish direct relationships between image features and crowd counts but struggled with generalization across different scenes and crowd densities. The emergence of density estimation methods marked a significant advancement, as they generate spatial density maps that provide not only total counts but also spatial distribution information of crowds.\\
% The advent of deep learning has revolutionized crowd counting research, with convolutional neural networks (CNNs) becoming the dominant approach for density map generation. In recent years, deep learning techniques \cite{wang2025comprehensive} have addressed many limitations of traditional methods, particularly in handling complex scenarios with varying illumination, severe occlusion, and perspective distortion. However, significant challenges remain, particularly in addressing scale variation, which represents one of the most persistent problems in crowd counting.\\
% The problem of scale variation is an ingrained and drastic challenge in crowd counting, and it severely degrades the performance of counting \cite{fan2022survey}. This challenge arises from the natural perspective effects in images, where individuals appear at vastly different scales depending on their distance from the camera. The traditional multi-scale neural networks exhibit larger scales for targets in the near range and become very small in the far range, which presents a significant challenge in meeting the task of crowd counting with significant scale differences \cite{zhou2025crowd}.\\
% Recent advances in crowd counting have focused on developing sophisticated architectures that can handle multi-scale variations while maintaining computational efficiency. The field of crowd counting is moving towards combining multiple methods and requires fresh, targeted datasets. Despite advancements, the field still faces challenges such as handling real-world scenarios and processing large crowds in real-time \cite{deng2024deep}. Modern approaches increasingly incorporate attention mechanisms, multi-scale feature extraction, and advanced fusion strategies to improve counting accuracy across diverse scenarios.\\
% In this context, the Focus and Fusion Network (FFNet) represents a significant contribution to addressing the scale variation challenge in crowd counting. FFNet introduces an innovative approach that combines focused attention mechanisms with multi-scale feature fusion to achieve superior performance in handling crowds with varying densities and scales. The network's architecture specifically addresses the limitations of existing methods by implementing targeted feature extraction and intelligent fusion strategies that enhance counting accuracy while maintaining computational efficiency.\\
% This work presents a comprehensive evaluation of FFNet within a deep stream system framework, demonstrating its effectiveness in real-world crowd counting scenarios and its potential for deployment in practical surveillance and monitoring applications. The integration of frequency exploitation techniques with crowd counting represents a novel approach that leverages both spatial and frequency domain information to achieve enhanced performance in challenging crowd scenarios.
\subsection{Motivations}
The development of an effective crowd counting system within a deep stream framework addresses critical limitations in current surveillance and monitoring solutions. Traditional crowd monitoring relies on manual observation or basic motion detection, which fails to provide accurate real-time population estimates in complex environments with high-density crowds, varying lighting, and diverse architectural settings.\\
The integration of crowd counting with super-resolution in a unified system addresses a fundamental gap where existing methods suffer from degraded performance on low-resolution streams—common in real-world deployments due to bandwidth and hardware constraints. This creates a need for systems that simultaneously enhance image quality and perform accurate crowd estimation.\\
The problem of scale variation is an ingrained and drastic challenge in crowd counting, and it severely degrades the performance of counting \cite{fan2022survey}. In practical scenarios, individuals appear at vastly different scales due to perspective effects and varying crowd densities, necessitating sophisticated architectures like FFNet that handle multi-scale variations while maintaining computational efficiency.\\
The increasing demand for intelligent crowd management in urban environments, public events, and transportation hubs drives the need for automated systems providing continuous, accurate crowd density information to support public safety decisions, resource allocation, and emergency response.
\subsection{Goals}
The primary goal is to develop a comprehensive deep stream system integrating frequency exploitation super-resolution with FFNet-based crowd counting to achieve superior performance in challenging real-world scenarios.\\
The first specific goal is implementing an effective solution to scale variation through FFNet's focus and fusion mechanisms. The system must handle crowds with significant scale differences, ensuring consistent accuracy regardless of perspective distortion or camera positioning through robust multi-scale feature extraction and intelligent fusion strategies.\\
The second goal is achieving real-time processing capabilities for continuous streaming applications. The system must maintain computational efficiency while delivering high accuracy, enabling deployment in resource-constrained environments by optimizing the balance between accuracy and processing speed.\\
The third goal focuses on enhancing robustness across diverse environmental conditions including different lighting, weather, architectural settings, and crowd behaviors. This includes handling severe occlusions, irregular crowd distributions, and dynamic lighting commonly encountered in practical applications.\\
Additionally, the system aims to provide comprehensive spatial crowd distribution through accurate density maps, enabling advanced analysis capabilities including crowd flow estimation, congestion detection, and anomaly identification for sophisticated crowd management applications.\\
Finally, the system establishes a framework for seamless integration with existing surveillance infrastructures through compatibility with standard video streaming protocols, facilitating easy deployment without extensive infrastructure modifications.